\chapter{Rastreabilidade Arquitetural e Atributos de Qualidade}
\label{cap_atributos_qualidade}

A validação de uma arquitetura de software fundamenta-se na sua capacidade comprovada de atender aos atributos de qualidade (\textit{Quality Attributes}) estipulados pelas partes interessadas e formalizados nos Requisitos Não Funcionais \cite{bass2021,richards2020}. Neste capítulo, estabelece-se o mapeamento de rastreabilidade entre as decisões tomadas nas cinco visões do Modelo 4+1 e os 10 Requisitos Não Funcionais (RNF01 a RNF10) do SIGEA-GTT, parametrizados pela norma internacional \textbf{ISO/IEC 25010} (SQuaRE) \cite{iso25010}.

% ===========================================================================
\section{Matriz de Rastreabilidade Arquitetural}
\label{sec_matriz_rastreabilidade_arq}
% ===========================================================================

A \autoref{tab_rastreabilidade_arq} correlaciona cada Requisito Não Funcional com sua respectiva dimensão de qualidade SQuaRE, as visões do modelo 4+1 impactadas e as táticas arquiteturais concretizadas na implementação.

\begin{table}[htb]
\centering
\small
\caption{Matriz de Rastreabilidade: Requisitos Não Funcionais vs. Modelo 4+1.}
\label{tab_rastreabilidade_arq}
\begin{tabularx}{\textwidth}{l >{\raggedright\arraybackslash}p{2.8cm} >{\raggedright\arraybackslash}p{2.8cm} X}
\toprule
\textbf{RNF} & \textbf{Categoria ISO 25010} & \textbf{Visões 4+1} & \textbf{Táticas e Decisões Arquiteturais Implementadas} \\
\midrule
\textbf{RNF01} & Usabilidade & Lógica, Física, Cenários & Layout projetado exclusivamente para desktops clínicos WUXGA, HD+, HD e Ultrawide 21:9 via TailwindCSS e componentes densos do Angular Material, sem foco mobile. \\
\midrule
\textbf{RNF02} & Desempenho & Lógica, Processo, Física & Latência da API REST $P_{95} \le 200\text{ ms}$; pool HikariCP, paginação \texttt{LIMIT/OFFSET} no banco, índices B-Tree e projeções MapStruct. \\
\midrule
\textbf{RNF03} & Segurança & Lógica, Processo, Desenvolvimento & Autenticação \textit{stateless} JWT HMAC-SHA256, BCrypt com fator de custo 12, filtros Spring Security 6 e guardas de rota Angular. \\
\midrule
\textbf{RNF04} & Conformidade / LGPD & Lógica, Física & Segregação física e de rede absoluta em relação aos bancos do HU-UFS/SUS; dados 100\% sintéticos simulados no PostgreSQL. \\
\midrule
\textbf{RNF05} & Usabilidade & Lógica, Desenvolvimento & Padrão universal de paginação em exatamente 10 registros por página gerenciado centralmente por componentes reutilizáveis. \\
\midrule
\textbf{RNF06} & Ergonomia / Design & Lógica, Cenários & Paleta hospitalar suave calibrada para baixo cansaço visual (azuis clínicos, verde esmeralda, cinzas neutros, âmbar e carmim cirúrgico). \\
\midrule
\textbf{RNF07} & Usabilidade & Lógica, Processo & Interceptor HTTP gerencia Signal \texttt{isLoading} para ativação em menos de 100 ms do \textit{spinner} com desfoque e Toasts flutuantes. \\
\midrule
\textbf{RNF08} & Confiabilidade & Lógica, Cenários & Diálogos modais em duas etapas com bloqueio de fluxo para ações destrutivas (exclusões e encerramentos de turma). \\
\midrule
\textbf{RNF09} & Confiabilidade & Desenvolvimento & \textit{Quality Gate} inegociável de 100\% de cobertura de ramos e linhas no backend (JUnit 5 + Jacoco) e frontend (Jest/Karma). \\
\midrule
\textbf{RNF10} & Manutenibilidade & Desenvolvimento & Documentação técnica integral da base de código gerada sem erros: Javadoc formal, SpringDoc OpenAPI 3 e Compodoc. \\
\bottomrule
\end{tabularx}
\fonte{Elaborado pelo autor (2026).}
\end{table}

% ===========================================================================
\section{Análise de Trade-offs Arquiteturais}
\label{sec_tradeoffs}
% ===========================================================================

Toda decisão arquitetural de engenharia envolve concessões conscientes (\textit{trade-offs}) entre atributos de qualidade concorrentes \cite{bass2021,richards2020}:

\begin{enumerate}
    \item \textbf{Flexibilidade Semiestruturada versus Integridade Relacional}:
    \begin{itemize}
        \item \textit{Decisão}: Adoção do modelo híbrido no PostgreSQL 16 (tabelas relacionais com chaves estrangeiras rígidas para entidades mestre e colunas nativas \texttt{JSONB} para Prontuários e Ferramentas da Qualidade);
        \item \textit{Benefício}: Permite que professores elaborem prontuários com seções médicas altamente dinâmicas sem necessidade de migrações DDL contínuas no banco;
        \item \textit{Compensação}: Validação da estrutura JSON delegada à camada de serviço e a validadores de esquema no backend Spring Boot.
    \end{itemize}

    \item \textbf{Segurança Criptográfica e Stateless versus Sobrecarga de Rede}:
    \begin{itemize}
        \item \textit{Decisão}: Autenticação puramente \textit{stateless} com JSON Web Tokens trafegados no cabeçalho \texttt{Authorization};
        \item \textit{Benefício}: Elimina o consumo de memória de sessão na JVM e permite escalabilidade horizontal irrestrita;
        \item \textit{Compensação}: Ligeiro acréscimo no tamanho de cada cabeçalho HTTP (aproximadamente 500 bytes por requisição), perfeitamente absorvido pela rede local de alta velocidade da UFS.
    \end{itemize}

    \item \textbf{Densidade Visual Clínica versus Suporte a Dispositivos Móveis}:
    \begin{itemize}
        \item \textit{Decisão}: Homologação estrita para resoluções desktop (WUXGA, HD+, HD e Ultrawide 21:9), sem suporte para smartphones;
        \item \textit{Benefício}: Permite a visualização simultânea de evoluções médicas, exames, cronômetro de 20 minutos e matrizes de causa-raiz sem sobrecarga cognitiva ou rolagem excessiva;
        \item \textit{Compensação}: O sistema exige o uso de estações de trabalho de laboratório ou computadores portáteis com telas de dimensões adequadas ao trabalho clínico hospitalar.
    \end{itemize}
\end{enumerate}
